\section*{Sets and Indices}

\begin{itemize}
    \item $I$: set of activities, with
    \[
    I=\{A,B,C,D,E,F,G\}.
    \]
    \item $P \subseteq I \times I$: set of precedence arcs $(i,j)$ such that activity $i$ must finish before activity $j$ can start.
\end{itemize}

\section*{Parameters}

\begin{itemize}
    \item $d_i$: duration of activity $i$ in days.
    \[
    d_A=4,\;
    d_B=3,\;
    d_C=5,\;
    d_D=2,\;
    d_E=10,\;
    d_F=10,\;
    d_G=1.
    \]
    \item $c^{\text{work}}$: project work cost per day,
    \[
    c^{\text{work}}=1000.
    \]
    \item $c^{\text{mach}}$: machine rental cost per day,
    \[
    c^{\text{mach}}=5000.
    \]
    \item Precedence set
    \[
    P=\{(A,G),(A,D),(E,F),(G,F),(D,C),(F,C),(F,B)\}.
    \]
\end{itemize}

\section*{Decision Variables}

\begin{itemize}
    \item $S_i$ (code: \texttt{S[a]}): start time of activity $i$, for all $i \in I$.
    \item $T$ (code: \texttt{T}): project time variable used in the objective and constrained to cover the completion of activities $B$ and $C$.
\end{itemize}

All variables are continuous and nonnegative:
\[
S_i \ge 0 \quad \forall i \in I, \qquad T \ge 0.
\]

\section*{Objective}

Minimize total cost:
\[
\min \; c^{\text{work}} T + c^{\text{mach}}\bigl(S_B + d_B - S_A\bigr).
\]

With the given data, this is
\[
\min \; 1000\,T + 5000\,(S_B + 3 - S_A).
\]

\section*{Constraints}

\begin{enumerate}
    \item \textbf{Precedence constraints.} For every $(i,j)\in P$,
    \[
    S_j \ge S_i + d_i.
    \]
    Explicitly,
    \begin{align*}
        S_G &\ge S_A + d_A,\\
        S_D &\ge S_A + d_A,\\
        S_F &\ge S_E + d_E,\\
        S_F &\ge S_G + d_G,\\
        S_C &\ge S_D + d_D,\\
        S_C &\ge S_F + d_F,\\
        S_B &\ge S_F + d_F.
    \end{align*}

    \item \textbf{Project time coverage.}
    \[
    T \ge S_C + d_C,
    \]
    \[
    T \ge S_B + d_B.
    \]
\end{enumerate}

\section*{Notes on Implementation}

The implemented model is a linear program solved by the code with a continuous LP solver interface. No integrality restrictions are imposed on start times.

The variable $T$ is not constrained to be the completion time of all activities; it is only required to be at least the completion times of activities $B$ and $C$. Given the precedence network, these are the only terminal activities explicitly used to define $T$ in the code.

The machine-rental term is modeled exactly as
\[
c^{\text{mach}}(S_B+d_B-S_A),
\]
i.e., the elapsed time from the start of activity $A$ to the finish of activity $B$. This is not generalized in the implementation to any broader machine-usage logic; the formulation above matches the code exactly.
